\section{The English of 2011}

\subsection{The regime under which the translation was made}

The English Order of Mass in use in the dioceses of the United States since Advent 2011 was made
under a stated and published rule, and the first task of an honest account is to state the rule
rather than to infer it from the result.

That rule is the instruction \emph{Liturgiam authenticam}, the fifth instruction for the right
implementation of the Constitution on the Sacred Liturgy, dated 28 March 2001 and published on 7
May 2001. It replaced the whole preceding body of translation norms: ``the norms set forth in
this Instruction are to be substituted for all norms previously published on the matter,'' with
one named exception; approvals already granted remain in effect even where they followed
different principles, but a new period begins for revising translations already
made.\endnote{\emph{Liturgiam authenticam} 8, 131, and 133, read 2026-07-25 in the official
English web state. The named exception at n.\ 8 is the 1994 instruction \emph{Varietates
legitimae}. The instruction's own explanation of why new norms were thought necessary is at nn.\
6--7: translations in various places were judged to need improvement, and the ``omissions or
errors which affect certain existing vernacular translations'' had impeded inculturation. This
study reports the instruction's account of the situation; it does not endorse or contest that
account. The earlier translation regime associated with the 1969 instruction on translation was
not examined for this study and is not characterised here.}

Its governing principle is stated at n.\ 20 and is worth having exactly:\endnote{\emph{Liturgiam authenticam} 20, read 2026-07-25. Quoted from the official English state
because the argument turns on this text's own wording; the promulgated Latin is in AAS 93 (2001)
685--726 and was not consulted.}

\begin{studyframe}
\emph{The translation of the liturgical texts of the Roman Liturgy is not so much a work of
creative innovation as it is of rendering the original texts faithfully and accurately into the
vernacular language. While it is permissible to arrange the wording, the syntax and the style in
such a way as to prepare a flowing vernacular text suitable to the rhythm of popular prayer, the
original text, insofar as possible, must be translated integrally and in the most exact manner,
without omissions or additions in terms of their content, and without paraphrases or glosses.}
\end{studyframe}

Four further norms shape the English Order of Mass specifically. The translation should be
easily understandable yet preserve the texts' dignity, beauty, and doctrinal precision (n.\ 25).
The liturgical texts are the voice of the Church at prayer rather than of particular
congregations, should be free of an overly servile adherence to prevailing modes of expression,
and may properly employ words that differ somewhat from everyday speech, since it is often by
this very fact that they become memorable and capable of expressing heavenly realities; the
instruction expects the gradual development of a recognisably sacred vernacular style (n.\ 27).
Varied Latin vocabulary should give rise to corresponding variety in the vernacular (n.\ 51).
And the Roman Rite's straightforward, concise, and compact manner of expression is to be
maintained insofar as possible (n.\ 57).\endnote{\emph{Liturgiam authenticam} 25, 27, 51, and
57, read 2026-07-25.}

Two provisions look directly at the cost of changing texts the people know. Without real
necessity, successive revisions should not notably change previously approved vernacular texts
of the Eucharistic Prayers which the faithful will have gradually committed to memory (n.\ 64);
and parts to be committed to memory, especially if sung, are to be changed only for a just and
considerable reason --- but where more significant changes are necessary to bring a text into
conformity with the instruction, it is preferable to make them at one time rather than
prolonging them over several editions, and a suitable period of catechesis should accompany the
publication (n.\ 74).\endnote{\emph{Liturgiam authenticam} 64 and 74, read 2026-07-25. Read
together, these two numbers describe exactly what happened in 2011: a single large change rather
than a series of small ones, accompanied by a programme of catechesis.} Whatever one thinks of
the 2011 English, the manner of its introduction --- everything at once, with preparation ---
was the manner the governing instrument prescribed.

\subsection{What actually changed, and on what authority}

The following are the changes in the people's parts of the Order of Mass that carry the most
weight, each with the instrument that accounts for it. The English wording is given only where
the argument turns on it, and never at a length that would reproduce the text.

{\footnotesize
\begin{longtable}{>{\raggedright\arraybackslash}p{0.20\linewidth}>{\raggedright\arraybackslash}p{0.30\linewidth}>{\raggedright\arraybackslash}p{0.44\linewidth}}
\toprule
\textbf{Place} & \textbf{Latin at issue} & \textbf{What accounts for the change}\\
\midrule
\endhead
Greeting response & \latin{Et cum spiritu tuo} & Named in \emph{Liturgiam authenticam} 56 as an expression of the ancient Church's heritage to be rendered as literally as possible. The 2011 English has ``and with your spirit.''\\
Penitential Act & \latin{mea culpa, mea culpa, mea maxima culpa} & Named in the same number of the same instruction. The triple form, absent from the 1973 English, is restored.\\
Creed, opening & \latin{Credo} & Singular in the Latin; n.\ 20's rule against additions and omissions in content. ``I believe'' replaces ``We believe.''\\
Creed & \latin{consubstantialem Patri} & Named in \emph{Liturgiam authenticam} 53 among Latin terms of rich meaning, for which transcription is one licensed solution. The 2011 English has ``consubstantial with the Father.'' The instruction licensed the solution; it did not mandate it.\\
Preface dialogue & \latin{Dignum et iustum est} & N.\ 20's integral rendering. ``It is right and just.''\\
\latin{Sanctus} & \latin{Dominus Deus Sabaoth} & N.\ 56's rule on ancient formulas; ``Lord God of hosts.''\\
Consecration of the chalice & \latin{qui pro vobis et pro multis effundetur} & The circular letter of 17 October 2006 (Prot.\ N.\ 467/05/L), which required a precise rendering in the next translation of the Roman Missal. ``For many'' replaces ``for all.''\\
Before Communion & \latin{ut intres sub tectum meum} & N.\ 43's rule that concrete biblical imagery keep its force by literal rendering; and \latin{IGMR} 84, which calls the response the prescribed words of the Gospel. ``Under my roof'' restores Matthew 8:8.\\
Orations & Latin periodic structure & Nn.\ 20, 51, and 57. The structural effect is described in this study without quotation, since the approved English of the orations is under copyright.\\
\bottomrule
\end{longtable}}

\subsection{\latin{Pro multis}: the one change with its own instrument}

The rendering of \latin{pro multis} is the only change in the 2011 English of the Order of Mass
that has a dedicated act of the Holy See behind it, and the act is worth setting out in full
because it is regularly misdescribed in both directions.

On 9 July 2005 the Congregation for Divine Worship, in agreement with the Congregation for the
Doctrine of the Faith, wrote to all presidents of episcopal conferences asking their opinion on
the translation of \latin{pro multis}. The replies were studied by both dicasteries and a report
was presented to Benedict XVI. On his directives, the Congregation wrote again on 17 October
2006, under the same protocol number, to all the presidents. The letter's Italian is the
original by its own preamble, and its substance is this.\endnote{Congregation for Divine Worship
and the Discipline of the Sacraments, circular letter Prot.\ N.\ 467/05/L of 17 October 2006, in
\emph{Notitiae} 42 (2006) 441--443, read publication-locally at artifact pages 57--59 of the
official Dicastery PDF on 2026-07-25. The Latin preamble printed above the letter records that
the Supreme Authority mandated its dispatch and that the Italian is to be held as the original
from which the other versions were made. The summary that follows renders the letter's four
numbered paragraphs; short phrases are quoted from the Italian.}

\begin{enumerate}
\item A text corresponding to \latin{pro multis} has been the formula in use in the Latin Roman
Rite since the first centuries. Over the preceding thirty years some approved vernacular texts
adopted the interpretative translation \emph{for all}, \latin{per tutti}, or the like.
\item There is no doubt about the validity of Masses celebrated with a duly approved formula
containing an expression equivalent to \emph{for all}, as the Congregation for the Doctrine of
the Faith had already declared in 1974; that formula would correspond to a correct
interpretation of the Lord's intention expressed in the text, and it is a dogma of faith that
Christ died on the cross for all men and women.
\item There are nevertheless many arguments in favour of a more precise translation. The
synoptic Gospels refer specifically to ``the many''; it would have been entirely possible to
have ``for all'' in the redaction of the Gospel texts, and the formula given for the institution
narrative is ``for many,'' faithfully so translated in most modern biblical versions. The Latin
Roman Rite has always said \latin{pro multis} and never \latin{pro omnibus} in the consecration
of the chalice. The anaphoras of the various Eastern Rites contain the equivalent formula in
their own languages. ``For many'' is a faithful translation of \latin{pro multis}, while ``for
all'' is rather an explanation belonging properly to catechesis. ``For many,'' while remaining
open to including every single human person, also reflects that this salvation is not
accomplished in a quasi-mechanical way without one's own willing or participation. And in
conformity with \emph{Liturgiam authenticam} one should strive to remain more faithful to the
Latin texts of the typical editions.
\item Conferences in countries where \emph{for all} or its equivalent was then in use were
asked to undertake the necessary catechesis of the faithful over the next one or two years, so
as to prepare them for the introduction of a precise vernacular translation of \latin{pro
multis} in the next translation of the Roman Missal that the bishops and the Holy See would
approve.
\end{enumerate}

Three misdescriptions are thereby excluded. It is not the case that the change was ICEL's
preference or an English-speaking peculiarity: it was required of every conference then using an
equivalent of ``for all,'' by an instrument of the Holy See acting on papal directives. It is
not the case that the change implies that Masses said with ``for all'' were invalid or that
Christ did not die for all: the letter denies both in terms, and the second is affirmed as a
dogma of faith in the same paragraph. And it is not the case that the change was made without
reasons: six are given, and they are of different weights --- the scriptural and the
Latin-textual arguments are strong, the argument from the Eastern anaphoras is corroborative, and
the argument at (e) about the manner of salvation is theological rather than philological and is
the one on which a reader may most reasonably differ.

\subsection{Assessing the translation}

Three judgements can be made from the evidence assembled here, and they should be kept apart.

\emph{On conformity.} The 2011 English of the Order of Mass conforms to the instrument under
which it was made. The specific changes most complained of --- ``and with your spirit,''
``consubstantial,'' ``for many,'' ``under my roof,'' the triple \latin{mea culpa} --- are each
either named in the governing instruction or required by a dedicated act. Whatever else may be
said, the translation did what it was told to do, and the complaint that it was arbitrary does
not survive contact with \emph{Liturgiam authenticam} 43, 53, and 56.

\emph{On the rule itself.} Whether the rule was the right rule is a separate question, and it
is not settled by the translation's conformity to it. The strongest case against
\emph{Liturgiam authenticam} is not that literal rendering is wrong but that the instruction's
own n.\ 25 --- easily understandable language preserving dignity, beauty, and doctrinal
precision --- sets a standard that a literal rendering of complex Latin periods can fail, and
that a rule which produces sentences the hearer cannot parse at speaking pace has defeated the
Council's own priority of intelligible participation. That is a real tension inside the
instruction, not merely a complaint from outside it.

\emph{On the execution.} Conformity to a rule and quality of execution are different again. This
study is not in a position to assess the 2011 English sentence by sentence, because doing so
responsibly would require quoting the text at length, which the rights position forbids, and
because a judgement of that kind properly belongs to reviewers competent in English prose and
in Latin, several of whom disagree. What can be said is that the two levels of criticism are
regularly conflated: an objection to a particular English sentence is not an objection to the
principle of close translation, and an objection to the principle is not answered by defending
a sentence.

\subsection{A note on what this study could not do}

The reader will have noticed that no English of the Order of Mass has been printed in these
pages. That is deliberate and it has a cost, which is best stated plainly rather than hidden.
The approved English of the people's parts, the orations, and the Eucharistic Prayers is under
copyright, and this study neither reproduces it nor paraphrases it closely enough to
reconstruct. In consequence, three things that a fuller study would do are missing here: a
side-by-side comparison of the 1973 and 2011 English of any complete text; an assessment of the
2011 English of the Roman Canon, which is the longest and most contested single stretch of the
translation; and any judgement about the singability of the sung parts. Where an argument turned
on a short English phrase, that phrase has been given. Where it would have required the text,
the argument has been left undeveloped and said so.
